\documentclass[a4paper,11pt]{article}
\usepackage[T1]{fontenc}
\usepackage[utf8]{inputenc}
\usepackage{lmodern}
\usepackage{amsthm,amsmath,amsfonts,amssymb}

\newtheorem{claim}{Claim}

\title{Graduate Algebra 1: midsem exam}
\author{Aditya Padiyar}
\date{October 2026}

\begin{document}

\maketitle

\section{Polynomials and field extensions}

\begin{claim}
  If $F \subseteq K$ is a field extension of degree 2, then it is normal
\end{claim}

\begin{proof}
  We have $K=F(\alpha)$. Let $p \in F[x]$ be the minimal polynomial of $\alpha$. Then $K$ is the splitting field of $p$ and hence is normal.
\end{proof}

\begin{claim}
  Let $q(x)=x^4-2x^2-1$
  \begin{enumerate}
    \item $q(x)$ is irreducible over $\mathbb{Q}$
    \item $q(x)$ does not split completely over $\mathbb{R}$
  \end{enumerate}
\end{claim}

\begin{proof}
  Consider the following:
  \begin{enumerate}
    \item We have $q(x+1)=x^4+4x^3+4x^2-2$. Now use Eisenstein's criterion with $p=2$ to show that $q(x+1)$ is irreducible. 
          Since the map $p(x) \to p(x+1)$ is an automorphism of the ring $Q[x]$, $q(x)$ is also irreducible.
    \item This is easy to see
  \end{enumerate}
\end{proof}

\section{Fields of prime characteristic}

\begin{claim}
  Suppose $F$ is finite and $a,b,c \in F^\times$. Consider the polynomials $p_1(x,y)=x^2+y^2-a$ and $p_2(x,y)=x^2+by^2-c$ in $F[x,y]$. Then in $F$, $p_1$ and $p_2$ have roots
\end{claim}

\begin{proof}
  Consider the following:
  \begin{itemize}
    \item If $F$ has characteristic $2$, by the surjectivity of the Frobenius map, we can find solutions of the form $(\alpha,0)$ to both equations
    \item Otherwise, use the fact that $|\{x^2:x \in F\}|=\frac{q+1}{2}$
  \end{itemize}
  
\end{proof}

\begin{claim}
  Suppose $F$ is a field of characteristic $p$ and $K$ is an algebraic closure. Let $f(x)=x^p-x-a \in F[x]$. Then we have:
  \begin{enumerate}
    \item $f(x)=(x+\alpha)(x+\alpha+1)\dots(x+\alpha+p-1)$ in $K[x]$, for some $\alpha \in K$
    \item If $\alpha \notin F$ then $f$ is irreducible over $F$
  \end{enumerate} 
\end{claim}

\begin{proof}
  Consider the following:
  \begin{enumerate}
    \item We have $f(x+1)=f(x)$
    \item If $(x+\alpha + c_1)\dots(x+\alpha + c_m) \in F[x]$, then $m\alpha + c_1 + \dots + c_m \in F$. This forces $m=p$ 
  \end{enumerate}
\end{proof}

\begin{claim}
  Suppose we have the following:
  \begin{itemize}
    \item $p$ is odd, $K_1=\mathbb{F}_p(t)$ and $K_2$ is an algebraic closure of $K_1$
    \item $f(x)=x^{2p}-t \in K_1[x]$
    \item $f(\theta)=0$ with $\theta \in K_2$
  \end{itemize}
  
  Then we have the following:
  \begin{enumerate}
    \item $f$ is irreducible in $K_1[x]$
    \item $f(x)=(x-\theta)^p(x+\theta)^p$ in $K_2[x]$
    \item $\theta$ is not purely inseparable over $K_1$
  \end{enumerate}
\end{claim}

\begin{proof}
  Consider the following:
  \begin{enumerate}
    \item $f$ is irreducible by Eisenstein's criterion
    \item Use the Frobenius homomorphism
    \item If $\theta^{p^n} = \theta^p t^{(p^{n-1}-1)} \in K_1$ then $\theta^p \in K_1$. But $\theta$ has degree $2p$ over $K_1$.
  \end{enumerate}
  
\end{proof}

\section{Modules}

Suppose we have the following:
\begin{itemize}
  \item $k$ is a field
  \item $I=t^2k[t]$
  \item $R=k+I$
  \item Regard $k[t]$ as an $R$ module using the inclusion $R \subseteq k[t]$
\end{itemize}

\begin{claim}
  The quotient vector space $I/I^2$ has basis $(t^2 + I^2,t^3 + I^2)$
\end{claim}

\begin{claim}
  $I$ is an ideal but not a principal ideal in $R$
\end{claim}

\begin{proof}
  Consider the following:
  \begin{itemize}
    \item It is easy to see that $I$ is an ideal of $R$
    \item Suppose $I$ is generated by a monic polynomial $p(t) \in I$. Then $p(t)$ is the lowest degree monic polynomial in $I$. 
    Hence $p(t)=t^2$. But the element $t^3 \in I$ does not belong to $t^2R$
  \end{itemize}
\end{proof}

\begin{claim}
  $k[t]$ is not a free $R$-module
\end{claim}

\begin{proof}
  Consider the following:
  \begin{itemize}
    \item The equation $(t^2p_2)p_1 - (t^2p_1)p_2=0$ shows that no two elements of $k[t]$ are linearly independent
    \item Suppose the module homomorphism $R \to k[t]$ given by $q(t) \to q(t)p(t)$ is surjective for some $p(t) \in k[t]$.
          We then get a surjective module homomorphism from $R/I$ to $k[t]/I$.
          Since $R/I$ is a 1 dimensional $k$ vector space we cannot have such a $k$ linear map
  \end{itemize}
\end{proof}

\begin{claim}
  There does not exist an $R$-module homomorphism $T:k[t] \to R$ such that $T(r)=r$ for all $r \in R$
\end{claim}

\begin{proof}
  If $T$ is such a map, we have $T(t^3)=t^3=t^2T(t)$ which means that $T(t)=t \notin R$ which is a contradiction
\end{proof}

\end{document}
